
\documentclass[UTF8]{amsart}
%Typical documenttypes: article/book
%some examples:
%\documentclass[reqno,11pt]{book}   %%%for books
%\documentclass[]{minimal}			%%%for Minimal Working Example


%for beamers, you have to change a lot. Especially, remove the package enumitem!!!



%%%%%%%%%%%%%%%%%%%% setting for fast compiling

%\special{dvipdfmx:config z 0}		% no compression

\includeonly{chapters/chapter9}		% In practice, use an empty document called "chapter9"	% usually for printing books






%%%%%%%%%%%%%%%%%%%% here we include packages

%%%basic packages for math articles
\usepackage{amssymb}
\usepackage{amsthm}
\usepackage{amsmath}
\usepackage{amsfonts}
\usepackage[shortlabels]{enumitem}	% It supersedes both enumerate and mdwlist. The package option shortlabels is included to configure the labels like in enumerate.

%%%packages for special symbols
\usepackage{pifont}					% Access to PostScript standard Symbol and Dingbats fonts
\usepackage{wasysym}				% additional characters
\usepackage{bm}						% bold fonts: \bm{...}
\usepackage{extarrows}				% may be replaced by tikz-cd
%\usepackage{unicode-math}			% unicode maths for math fonts, now I don't know how to include it
%\usepackage{ctex}					% Chinese characters, huge difference.


%%%basic packages for fancy electronic documents
\usepackage[colorlinks]{hyperref}
\usepackage[table,hyperref]{xcolor} 			% before tikz-cd. 
%\usepackage[table,hyperref,monochrome]{xcolor}	% disable colored output (black and white)

%%%packages for figures and tables (general setting)
\usepackage{float}				%Improved interface for floating objects
\usepackage{caption,subcaption}
\usepackage{adjustbox}			% for me it is usually used in tables 
\usepackage{stackengine}		%baseline changes

%%%packages for commutative diagrams
\usepackage{tikz-cd}
\usepackage{quiver}			% see https://q.uiver.app/

%%%packages for pictures
\usepackage[width=0.5,tiewidth=0.7]{strands}
\usepackage{graphicx}			% Enhanced support for graphics

%%%packages for tables and general settings
\usepackage{array}
\usepackage{makecell}
\usepackage{multicol}
\usepackage{multirow}
\usepackage{diagbox}
\usepackage{longtable}

%%%packages for ToC, LoF and LoT







 %https://tex.stackexchange.com/questions/58852/possible-incompatibility-with-enumitem










%%%%%%%%%%%%%%%%%%%% here we include theoremstyles

\numberwithin{equation}{section}

\theoremstyle{plain}
\newtheorem{theorem}{Theorem}[section]

\newtheorem{setting}[theorem]{Setting}
\newtheorem{definition}[theorem]{Definition}
\newtheorem{lemma}[theorem]{Lemma}
\newtheorem{proposition}[theorem]{Proposition}
\newtheorem{corollary}[theorem]{Corollary}
\newtheorem{conjecture}[theorem]{Conjecture}

\newtheorem{claim}[theorem]{Claim}
\newtheorem{eg}[theorem]{Example}
\newtheorem{ex}[theorem]{Exercise}
\newtheorem{fact}[theorem]{Fact}
\newtheorem{ques}[theorem]{Question}
\newtheorem{warning}[theorem]{Warning}



\newtheorem*{bbox}{Black box}
\newtheorem*{notation}{Conventions and Notations}


\numberwithin{equation}{section}


\theoremstyle{remark}

\newtheorem{remark}[theorem]{Remark}
\newtheorem*{remarks}{Remarks}

%%% for important theorems
%\newtheoremstyle{theoremletter}{4mm}{1mm}{\itshape}{ }{\bfseries}{}{ }{}
%\theoremstyle{theoremletter}
%\newtheorem{theoremA}{Theorem}
%\renewcommand{\thetheoremA}{A}
%\newtheorem{theoremB}{Theorem}
%\renewcommand{\thetheoremB}{B}







%%%%%%%%%%%%%%%%%%%% here we declare some symbols

%%%%%%%DeclareMathOperator
%see here for why newcommand is better for DeclareMathOperator: https://tex.stackexchange.com/questions/67506/newcommand-vs-declaremathoperator

%%%%%basic symbols. Keep them!

%%%symbols for sets and maps
\DeclareMathOperator{\pt}{\operatorname{pt}}	%points. Other possibilities are \{pt\}, \{*\}, pt, * ...
\DeclareMathOperator{\Id}{\operatorname{Id}}	%identity in groups.
\DeclareMathOperator{\Img}{\operatorname{Im}}

\DeclareMathOperator{\Ob}{\operatorname{Ob}}
\DeclareMathOperator{\Mor}{\operatorname{Mor}}	%difference of Mor and Hom: Hom is usually for abelian categories
\DeclareMathOperator{\Hom}{\operatorname{Hom}}	\DeclareMathOperator{\End}{\operatorname{End}}
\DeclareMathOperator{\Aut}{\operatorname{Aut}}

%%%symbols for linear algebras and 
%%linear algebras
\DeclareMathOperator{\tr}{\operatorname{tr}}
\DeclareMathOperator{\diag}{\operatorname{diag}}	%for diagonal matrices

%%abstract algebras
\DeclareMathOperator{\ord}{\operatorname{ord}}
\DeclareMathOperator{\gr}{\operatorname{gr}}
\DeclareMathOperator{\Frac}{\operatorname{Frac}}

%%%symbols for basic geometries
\DeclareMathOperator{\vol}{\operatorname{vol}}	%volume
\DeclareMathOperator{\dist}{\operatorname{dist}}
\DeclareMathOperator{\supp}{\operatorname{supp}}

%%%symbols for category
%%names of categories
\DeclareMathOperator{\Mod}{\operatorname{Mod}}
\DeclareMathOperator{\Vect}{\operatorname{Vect}}


%%%symbols for homological algebras
\DeclareMathOperator{\Tor}{\operatorname{Tor}}
\DeclareMathOperator{\Ext}{\operatorname{Ext}}
\DeclareMathOperator{\gldim}{\operatorname{gl.dim}}
\DeclareMathOperator{\projdim}{\operatorname{proj.dim}}
\DeclareMathOperator{\injdim}{\operatorname{inj.dim}}
\DeclareMathOperator{\rad}{\operatorname{rad}}


%%%symbols for algebraic groups
\DeclareMathOperator{\GL}{\operatorname{GL}}
\DeclareMathOperator{\SL}{\operatorname{SL}}

%%%symbols for typical varieties
\DeclareMathOperator{\Gr}{\operatorname{Gr}}
\DeclareMathOperator{\Flag}{\operatorname{Flag}}

%%%symbols for basic algebraic geometry
\DeclareMathOperator{\Spec}{\operatorname{Spec}}
\DeclareMathOperator{\Coh}{\operatorname{Coh}}
\newcommand{\Dcoh}{\mathcal{D}_{\operatorname{Coh}}}%%%This one shows the difference between \DeclareMathOperator and \newcommand
\DeclareMathOperator{\Pic}{\operatorname{Pic}}
\DeclareMathOperator{\Jac}{\operatorname{Jac}}

%%%%%advanced symbols. Choose the part you need!

%%%symbols for algebraic representation theory
\DeclareMathOperator{\ind}{\operatorname{ind}}	%\ind(Q) means the set of  equivalence classes of finite dimensional indecomposable representations
\DeclareMathOperator{\Res}{\operatorname{Res}}
\DeclareMathOperator{\Ind}{\operatorname{Ind}}
\DeclareMathOperator{\cInd}{\operatorname{c-Ind}}


\DeclareMathOperator{\Rep}{\operatorname{Rep}}
\DeclareMathOperator{\rep}{\operatorname{rep}} %usually rep means the category of finite dimensional representations, while Rep means the category of representations.
\DeclareMathOperator{\Irr}{\operatorname{Irr}}
\DeclareMathOperator{\irr}{\operatorname{irr}}
\DeclareMathOperator{\Adm}{\operatorname{\Pi}}
\DeclareMathOperator{\Char}{\operatorname{Char}}
\DeclareMathOperator{\WDrep}{\operatorname{WDrep}}

%%%symbols for algebraic topology
\DeclareMathOperator{\EGG}{\operatorname{E}\!}
\DeclareMathOperator{\BGG}{\operatorname{B}\!}

\DeclareMathOperator{\chern}{\operatorname{ch}^{*}}
\DeclareMathOperator{\Td}{\operatorname{Td}}
\DeclareMathOperator{\AS}{\operatorname{AS}}	%Atiyah--Segal completion theorem 

%%%symbols for Auslander--Reiten theory 
\DeclareMathOperator{\Modup}{\overline{\operatorname{mod}}}
\DeclareMathOperator{\Moddown}{\underline{\operatorname{mod}}}
\DeclareMathOperator{\Homup}{\overline{\operatorname{Hom}}}
\DeclareMathOperator{\Homdown}{\underline{\operatorname{Hom}}}


%%%symbols for operad
\DeclareMathOperator{\Com}{\operatorname{\mathcal{C}om}}
\DeclareMathOperator{\Ass}{\operatorname{\mathcal{A}ss}}
\DeclareMathOperator{\Lie}{\operatorname{\mathcal{L}ie}}
\DeclareMathOperator{\calEnd}{\operatorname{\mathcal{E}nd}} %cal=\mathcal


%%%%%personal symbols. Use at your own risk!

%%%symbols only for master thesis
\DeclareMathOperator{\ptt}{\operatorname{par}}	%the partition map
\DeclareMathOperator{\str}{\operatorname{str}}	%strict case
\DeclareMathOperator{\RRep}{\widetilde{\operatorname{Rep}}}
\DeclareMathOperator{\Rpt}{\operatorname{R}}
\DeclareMathOperator{\Rptc}{\operatorname{\mathcal{R}}}
\DeclareMathOperator{\Spt}{\operatorname{S}}
\DeclareMathOperator{\Sptc}{\operatorname{\mathcal{S}}}
\DeclareMathOperator{\Kcurl}{\operatorname{\mathcal{K}}}
\DeclareMathOperator{\Hcurl}{\operatorname{\mathcal{H}}}
\DeclareMathOperator{\eu}{\operatorname{eu}}
\DeclareMathOperator{\Eu}{\operatorname{Eu}}
\DeclareMathOperator{\dimv}{\operatorname{\underline{\mathbf{dim}}}}
\DeclareMathOperator{\St}{\mathcal{Z}}

%%%%%symbols which haven't been classified. Add your own math operators here!


\DeclareMathOperator{\Modr}{\operatorname{-Mod}}
\DeclareMathOperator{\Sym}{\operatorname{Sym}}
\DeclareMathOperator{\rk}{\operatorname{rk}}
\DeclareMathOperator{\Hnm}{\operatorname{H}}


%%%%%%%newcommand

%%%basic symbols
\newcommand{\norm}[1]{\Vert{#1}\Vert}

%%%symbols only for master thesis
\newcommand{\dimvec}[1]{\mathbf{#1}}
\newcommand{\abdimvec}[1]{|\dimvec{#1}|}
\newcommand{\ftdimvec}[1]{\underline{\dimvec{#1}}}

\newcommand{\absgp}[1]{\mathbb{#1}}
\newcommand{\WWd}{\absgp{W}_{\abdimvec{d}}}
\newcommand{\Wd}{W_{\dimvec{d}}}
\newcommand{\MinWd}{\operatorname{Min}(\absgp{W}_{\abdimvec{d}},W_{\dimvec{d}})}
\newcommand{\Compd}{\operatorname{Comp}_{\dimvec{d}}}
\newcommand{\Shuffled}{\operatorname{Shuffle}_{\dimvec{d}}}

\newcommand{\Omcell}{\Omega}
\newcommand{\OOmcell}{\boldsymbol{\Omega}}
\newcommand{\Vcell}{\mathcal{V}}
\newcommand{\VVcell}{\boldsymbol{\mathcal{V}}}
\newcommand{\Ocell}{\mathcal{O}}
\newcommand{\OOcell}{\boldsymbol{\mathcal{O}}}
\newcommand{\preimage}[1]{\widetilde{#1}}
\newcommand{\orde}{\operatorname{ord}_e}
\newcommand{\fakestar}{*}

%as the subscription of Hom
\newcommand{\Alggp}{\text{-Alg gp}}

%as automorphic representations
\newcommand{\Acusp}{\mathcal{A}_{\operatorname{cusp}}}




%%%%%%%%%%%%%%%%%%%% here we make some blocks for special features. 

%%%% todo notes %%%%
\usepackage[colorinlistoftodos,textsize=footnotesize]{todonotes}
\setlength{\marginparwidth}{2.5cm}
\newcommand{\leftnote}[1]{\reversemarginpar\marginnote{\footnotesize #1}}
\newcommand{\rightnote}[1]{\normalmarginpar\marginnote{\footnotesize #1}\reversemarginpar}









%%%%%%%%%%%%%%%%%%%% here we make some global settings. Understand everything here before you make a document!

\usepackage[a4paper,left=3cm,right=3cm,bottom=4cm]{geometry}
\usepackage{indentfirst}	% Indent first paragraph after section header

\setcounter{tocdepth}{2}


%https://latexref.xyz/_005cparindent-_0026-_005cparskip.html
\setlength{\parindent}{15pt}	
\setlength{\parskip}{0pt plus1pt}

%\setlength\intextsep{0cm}
%\setlength\textfloatsep{0cm}
\def\arraystretch{1}
%\setcounter{secnumdepth}{3}

\allowdisplaybreaks


\begin{document}

% The beginning depends on the documentclass. Rewrite this part if you use different documentclass!
\date{\today}

\title
{A counterexample to generic injectivity of the tangent Gauss map
}
\author{Xiaoxiang Zhou}
\address{Institut für Mathematik\\
Humboldt-Universität zu Berlin\\
Berlin, 12489\\ Germany\\} 
\email{email:xiaoxiang.zhou@hu-berlin.de}


\maketitle
\tableofcontents

\begin{theorem}
There exists a nondegenerate curve $C$ on a simple abelian variety $A$ of dimension $n>2$ such that $\phi_C$ is not generically injective and no nontrivial automorphism $\tau \in \Aut(A)$ preserves $C$ and acts as a scalar on $T_0A$.
\end{theorem}
\section{Strategy}

Let $C$ be a nonhyperelliptic curve, and let $\mathcal{L}$ be a line bundle of degree $3$ with $
h^0(\mathcal{L})=2$. Then $\mathcal{L}$ is base-point-free and induces a map $|\mathcal{L}|: C \longrightarrow \mathbb{P}^1$ of degree $3$.

The linear system $\left(\mathcal{L}^{\otimes 2}, \Sym^2 \Hnm^0(\mathcal{L})\right)$ induces a map $\phi:C \longrightarrow \mathbb{P}^2$ that factors through $|\mathcal{L}|$. Thus $\phi$ is not generically injective. Our aim is to realize $\phi$ as the tangent Gauss map of a curve $C \subset A$. This requires:

\begin{itemize}
\item $\mathcal{L}^{\otimes 2}= \omega_C$;
\item $\Sym^2 \Hnm^0(\mathcal{L}) \subset \Hnm^0(\omega_C)$ is the cotangent space of an abelian subvariety.
\end{itemize}
We refer to the second requirement as the \emph{cotangent-space condition}.

If we first drop the nonhyperellipticity assumption, we can find a curve satisfying these requirements. Take a hyperelliptic curve $C$ that admits a tangential cover $f:C \longrightarrow E$ at a Weierstrass point $p_0$. This means that the subspace $\Sym^2 \Hnm^0(\omega_C(-2p_0)) \subset \Hnm^0(\omega_C)$ satisfies the cotangent-space condition. For $g=4$, set $\mathcal{L}=\mathcal{O}_C(3p_0)$. Then
\begin{itemize}
\item $\mathcal{L}^{\otimes 2}= \mathcal{O}_C(6p_0) \cong \omega_C$;
\item $\Sym^2 \Hnm^0(\mathcal{L}) \cong \Hnm^0\left(\mathcal{O}_C(4p_0)\right) = \Hnm^0\left(\omega_C(-2p_0)\right)$ satisfies the cotangent-space condition.
\end{itemize}
We use this hyperelliptic curve as a seed and deform it to obtain a nonhyperelliptic curve while preserving the two requirements above.

\section{A dimension count for deformations}

We seek a local deformation $(C_t,\mathcal{L}_t)$ of the initial spin curve for which $\Sym^2 \Hnm^0(\mathcal{L}_t) \subset \Hnm^0(\omega_{C_t})$ continues to satisfy the cotangent-space condition and $C_t$ is nonhyperelliptic.

Let $B$ denote the local deformation space of the spin curve $(C_0,\mathcal{L}_0)$. Its dimension is
$$\dim B = h^1(\mathcal{T}_C)=3g-3=9.$$
\begin{enumerate}
\item The locus $\Theta_{null}= \left\{ (C,\mathcal{L}) \;\middle|\; h^0(\mathcal{L}) \geqslant 2 \right\}$ is a divisor.
\item Define the rational subspace
$$V_0:= H^1(C_0;\mathbb{Q}) \cap \Hnm^0\left(\omega_C(-2p_0)\right) \subset H^1(C_0;\mathbb{Q}),$$
and let $V_t \subset H^1(C_t;\mathbb{Q})$ be its parallel transport under the Gauss--Manin connection.

We also define the complex subspace 
$$W_0:=\Sym^2 \Hnm^0(\mathcal{L_0}) \subset \Hnm^0(\omega_{C_0}),$$
and let $W_t:=\Sym^2 \Hnm^0(\mathcal{L_t})\subset \Hnm^0(\omega_{C_t})$ be the corresponding subspace obtained by deforming the line bundle.

To ensure that $W_t$ continues to satisfy the cotangent-space condition, we use the following equivalent formulations:
\begin{equation*}
\begin{aligned}
   & \begin{cases}
   V_t \text{ remains a Hodge structure}\\
   V_{t,\mathbb{C}} \cap \Hnm^0(\omega_C)=W_t
   \end{cases} \\ 
  \Longleftrightarrow\;& \;V_{t,\mathbb{C}} = W_t \oplus \overline{W}_t  \\ 
  \Longleftrightarrow\;& \; W_t \subseteq V_{t,\mathbb{C}}\\
  \Longleftrightarrow\;& \; \rk \left[W_t \longrightarrow H^1(C_t;\mathbb{C})/ V_{t,\mathbb{C}}\right] =0
\end{aligned}
\end{equation*}
The last condition cuts out a locus of codimension $6$. Together with the divisor condition above, this leaves a deformation space of dimension $9-1-6=2$.

\item To find a nonhyperelliptic deformation, we compare this dimension with that of the local hyperelliptic branch. On a hyperelliptic curve $C_t$, a line bundle with $\deg \mathcal{L}_t=3$ and $h^0 (\mathcal{L}_t)=2$ must satisfy $\mathcal{L}_t \cong \mathcal{O}_{C_t}(3p_0)$ for some Weierstrass point $p_0$. Consequently,
$$W_t= \Sym^2 \Hnm^0(\mathcal{L}_t) = \Hnm^0\left(\omega_{C_t}(-2p_0)\right).$$
The following two lemmas show that the local hyperelliptic branch has dimension $1$:
\begin{lemma}
Let $C$ be a hyperelliptic curve and $p_0$ a Weierstrass point. The subspace
$\Hnm^0\left(\omega_{C_t}(-2p_0)\right)$ satisfies the cotangent-space condition if and only if there exists a tangential cover $f:C \longrightarrow E$ at $p_0$.
\end{lemma}
\begin{lemma}\cite[Proposition 2.5, 3]{treibich2025}
For a fixed elliptic curve $E$ and an integer $d$, there are only finitely many tangential covers $f:C \longrightarrow E$ of degree $d$.
\end{lemma}
\end{enumerate}

Thus, within a $3$-dimensional deformation space, the hyperelliptic locus has dimension at most $2$. It follows that there exists a nonhyperelliptic curve satisfying the required conditions.

%\nocite{Eberhardt2022Koszul}	% cite articles which are not cited in the document yet

% Remember to protect the uppercase of people's name and LaTeX symbols

\bibliographystyle{plain}
\bibliography{tangent_counterexample_references}
\end{document}
